\chapter{Розмова через повітря}
% full version: chapters/05_serocomputer.tex

\pencilfig{5}{\pencilart{5}}{Ліворуч телефон: провід від одного до одного. Праворуч радіо: одне повідомлення одразу всім навколо.}

Досі нейрони говорили одне з одним телефоном: провід до проводу, адресно. Тепер увімкнімо першу радіостанцію --- \nt{SERO}-нейрони, що викидають серотонін. Їх лише кілька сотень тисяч на весь мозок, але їхні проводи розходяться величезними його ділянками.

\section*{Чотири відмінності}

\nt{SERO}-нейрон не схожий на \nt{GLUT}. По-перше, знак його сигналу обирає отримувач: у серотоніну є «замки» і того, і іншого знака. По-друге, він працює сам: у стані неспання рівно клацає кілька разів на секунду, якщо отримує сталу фонову підживку, --- це метроном. По-третє, він повільний: його сигнал діє не за мілісекунди, а за сотні мілісекунд. По-четверте, він часто викидає речовину не у вузьку щілину, а в простір навколо, де вона розпливається, як крапля чорнила у воді. Отримувач відчуває концентрацію і не знає, від кого прийшла молекула.

\section*{Логіка, якої немає}

Метроном із гальмівним «замком» --- це готове «ні»: поки серотоніну немає, нейрон працює; з'явився --- замовк. Один нейрон замість цілої схеми. Здається, серотонінова мережа логічно сильніша за \nt{GLUT}-мережу.

Але є заковика. Речовина, викинута в простір, змішується. Два сигнали лишаються різними, лише якщо їх випустили в місцях, рознесених далі, ніж речовина встигає розпливтися, --- а це десятки мікрометрів. А в кожного серотонінового нейрона величезна кількість місць викиду, розкиданих великими ділянками, і «проводи» різних нейронів перекриваються. Окремих бітів такою мережею не передаси. Вона може передавати лише кілька загальних величин --- як радіо, яким усі слухають ту саму станцію.

\pencilfig{123}{%
% two release sites: far apart the clouds stay separate (two signals), close together they merge (one level)
\foreach \x/\y in {0.9/1.55, 4.1/1.55, 1.9/0, 2.9/0}{
  \foreach \r/\o in {0.95/0.07, 0.7/0.09, 0.45/0.12}{\fill[pcViolet, opacity=\o] (\x,\y) circle (\r);}
  \draw[dotpen=pcViolet, fill=pcViolet] (\x,\y) circle (0.07);}
\draw[arr] (1.95,1.55) -- (3.05,1.55); \draw[arr] (3.05,1.55) -- (1.95,1.55);
\node[lbl, anchor=south, align=center] at (2.5,1.6) {далі, ніж\\розпливається};
\node[lbl, anchor=west] at (5.25,1.55) {два сигнали};
\node[lbl, anchor=west] at (5.25,0) {одна спільна хмара};
}{Серотонін викидається в простір і розпливається хмарою. Два джерела дають два різні сигнали, лише якщо вони далі одне від одного, ніж речовина встигає розпливтися. Якщо ближче, хмари зливаються в один спільний рівень.}

\section*{Що вона робить насправді}

Для однієї загальної величини в цієї системи все є. Концентрація серотоніну згладжує окремі клацання в плавний рівень --- так найпростіший перетворювач перетворює потік імпульсів на напругу. Цей рівень тримається секунди й сам тане: не біт, який треба стирати, а слід, що повільно зникає.

Ще в серотонінових нейронів є «замки» на самих собі: викинутий серотонін пригальмовує власне джерело. Інженер упізнає в цьому підсилювач із від'ємним зворотним зв'язком --- термостат, що тримає рівень. Щоправда, досліди показують, що це самогальмування влаштоване адресно й спричиняє лише короткі паузи, тож роль термостата поки радше гіпотеза, ніж факт.

\section*{Скільки коштує чекати}

Яку загальну величину варто тримати повільною й однаковою для всієї мережі? Добрий кандидат --- \emph{терпіння}. Уявіть вибір: маленька винагорода зараз чи велика, але за кілька секунд. Скільки ви готові чекати, залежить від того, як швидко для вас «знецінюється» майбутнє. За однією з гіпотез, саме цю ручку й крутить серотонін. Є досліди на її користь: якщо штучно ввімкнути серотонінові нейрони, миша довше чекає відкладеної винагороди, перш ніж здатися.

Цей розділ вводить три пристрої, якими користуватимуться всі радіостанції мозку: нейрон-метроном, «провід», який насправді є об'ємом, і повільний рівень замість окремих клацань.

\fullversion{5}
